\section{Numerical experiments}
\label{sec:numerical}

All experiments use the implementation of Section~\ref{sec:method} with the fixed Newton tolerance
$\tau_N=10^{-11}$ and the dense automatic-differentiation Jacobian; the scripts that produce every
number and figure of this section are listed in \ref{app:repro}. Errors are measured
against a reference solution computed with the same method at a step size $h_{\rm ref}$ much
smaller than the finest step of the sweep; the difference between the references at $h_{\rm ref}$
and $2h_{\rm ref}$ is reported as the resolution floor of the reference, and a fitted order uses only
the step sizes whose errors exceed that floor by a factor of $100$. Position and velocity errors are
Frobenius norms over the bodies at the final time; the orientation error is the largest geodesic
angle between the numerical and reference rotations, and the angular-velocity error the Frobenius
norm of the body-frame difference.

\subsection{Convergence on the frictional chain}
\label{sec:e1}

The chain of Table~\ref{tab:chain} with $v_s=0.5$ is integrated over its regular branch $[0,0.5]$
with $h=0.1/2^k$, $k=0,\dots,5$, against the reference at $h_{\rm ref}=0.1/256$ ($1280$ steps).
Table~\ref{tab:e1} and Figure~\ref{fig:e1} show the result. All four error measures decrease by
$5.5$ to $6$ orders of magnitude over the sweep; the fitted orders over the five step pairs above
the floor are $6.25$ (position), $6.28$ (orientation), $6.11$ (velocity) and $6.27$ (angular
velocity), and the consecutive pairwise orders settle at $5.98$ and $5.99$ on the two finest
pairs. The endpoint velocity-level constraint norm is at roundoff on every step because of the
closure; the position-level norm, which is not projected, decreases as $h^6$ from $10^{-7}$ to
$3\times10^{-15}$ as the theory predicts. The Newton solve takes $6.0$ iterations per step on the
coarse grids and $5.7$ on the fine ones.

\begin{table}[!htb]
\centering
\caption{Convergence on the frictional chain, $T=0.5$, $v_s=0.5$: errors at $t=T$ against the
$h_{\rm ref}=0.1/256$ reference, pairwise order of the position error to the next finer step, and
maximal endpoint constraint norms at position ($\Phi$) and velocity ($\Phi_qv$) level; the number of steps is $0.5/h$. Resolution
floor of the reference: $3\times10^{-15}$ (position), $3\times10^{-15}$ (orientation),
$2\times10^{-14}$ (velocity), $1\times10^{-14}$ (angular velocity).}
\label{tab:e1}
\scriptsize
\setlength{\tabcolsep}{2.5pt}
\begin{tabular}{@{}lrllllrll@{}}
\toprule
$h$ & Newton/step & position & orientation & velocity & ang.\ velocity & order & $\|\Phi\|$ & $\|\Phi_qv\|$ \\
\midrule
$0.1$      & 6.00 & $2.77\mathrm{e}{-4}$  & $8.77\mathrm{e}{-5}$  & $4.96\mathrm{e}{-4}$  & $1.10\mathrm{e}{-3}$  & 6.95 & $9.8\mathrm{e}{-8}$  & $7.7\mathrm{e}{-16}$ \\
$0.05$     & 6.00 & $2.24\mathrm{e}{-6}$  & $7.98\mathrm{e}{-7}$  & $2.18\mathrm{e}{-6}$  & $9.47\mathrm{e}{-6}$  & 4.58 & $1.5\mathrm{e}{-9}$  & $1.3\mathrm{e}{-15}$ \\
$0.025$    & 5.90 & $9.38\mathrm{e}{-8}$  & $2.82\mathrm{e}{-8}$  & $1.95\mathrm{e}{-7}$  & $3.59\mathrm{e}{-7}$  & 7.64 & $2.4\mathrm{e}{-11}$ & $1.2\mathrm{e}{-15}$ \\
$0.0125$   & 5.72 & $4.69\mathrm{e}{-10}$ & $1.43\mathrm{e}{-10}$ & $9.42\mathrm{e}{-10}$ & $1.82\mathrm{e}{-9}$  & 5.98 & $3.7\mathrm{e}{-13}$ & $9.8\mathrm{e}{-14}$ \\
$0.00625$  & 5.72 & $7.45\mathrm{e}{-12}$ & $2.28\mathrm{e}{-12}$ & $1.49\mathrm{e}{-11}$ & $2.89\mathrm{e}{-11}$ & 5.99 & $6.3\mathrm{e}{-15}$ & $9.6\mathrm{e}{-14}$ \\
$0.003125$ & 5.72 & $1.17\mathrm{e}{-13}$ & $3.55\mathrm{e}{-14}$ & $2.33\mathrm{e}{-13}$ & $4.44\mathrm{e}{-13}$ & --   & $3.0\mathrm{e}{-15}$ & $8.2\mathrm{e}{-15}$ \\
\midrule
fitted order & & 6.25 & 6.28 & 6.11 & 6.27 & & & \\
\bottomrule
\end{tabular}
\end{table}

\paperfig{\figpath/e1_convergence.png}
{Convergence on the frictional chain over the regular branch $[0,0.5]$: errors at $t=T$ (left) and
over all grid points (right) against the $h_{\rm ref}=0.1/256$ reference, with the slope-6 line
anchored at the finest position error and the resolution floor of the reference.}
{fig:e1}

\paragraph{Exact stage identity and constants of the analysis}
Three diagnostics of Section~\ref{sec:order} are evaluated on the chain runs.
(i)~Each converged stage vector, obtained with a Newton tolerance of $10^{-13}$, is read in the
joint-coordinate chart and the reduced collocation defects $\Delta_i[s_j]$, $\Delta_i[\dot s_j]$,
$\Delta_i[\theta_j]$, $\Delta_i[\dot\theta_j]$ of Eq.~\eqref{eq:collocation-defect-operator} are
evaluated together with the component of $d_j,\dot d_j,\ddot d_j$ orthogonal to $n_j$: all are at
roundoff (Table~\ref{tab:exact-identity-check}), as Lemma~\ref{lem:exact-stage-identity} states.
(ii)~The constants of Lemma~\ref{lem:p2-from-p1} are measured from the implemented Jacobian at the
converged stage and at the root of the $h=0$ stage system (Table~\ref{tab:p2-constants-check}): the
conclusion $\|J_h^{-1}\|\le2\|J_0^{-1}\|$ holds with ratio at most $1.55$, while the sufficient
Neumann condition would need $h\le2\times10^{-5}$ in the Euclidean norm and $h\le8\times10^{-4}$ in
the endpoint-linearized norm; without joint friction the threshold rises to $1.4\times10^{-2}$.
(iii)~The distances of the two predictors to $Z_G$ (Table~\ref{tab:predictor-check}) confirm the
discussion after Lemma~\ref{lem:newton-envelope}: the lifted endpoint predictor is $O(h)$ from the
Gauss stage, the implemented predictor is at distance about $50$ and its first Newton step expands
that distance by up to $18\%$.

\begin{table}[!htb]
\centering
\caption{Exact stage identity at the converged stages of the chain (Newton tolerance $10^{-13}$,
window $[0,0.08]$): maximal implemented residual, maximal reduced collocation defect and maximal
component of $d_j,\dot d_j,\ddot d_j$ orthogonal to $n_j$.}
\label{tab:exact-identity-check}
\small
\begin{tabular}{P{0.07\linewidth}P{0.08\linewidth}P{0.13\linewidth}P{0.17\linewidth}P{0.17\linewidth}P{0.17\linewidth}}
\toprule
$h$ & steps & Newton total & max stage residual & max reduced defect & max $\perp n_j$ component \\
\midrule
0.04 & 2 & 12 & $6.8\times10^{-14}$ & $9.0\times10^{-15}$ & $1.5\times10^{-15}$ \\
0.02 & 4 & 24 & $1.5\times10^{-14}$ & $5.5\times10^{-16}$ & $1.9\times10^{-15}$ \\
0.01 & 8 & 48 & $1.6\times10^{-14}$ & $5.8\times10^{-16}$ & $2.0\times10^{-15}$ \\
\bottomrule
\end{tabular}
\end{table}

\begin{table}[!htb]
\centering
\caption{Constants of Lemma~\ref{lem:p2-from-p1} on the chain: $M_0=\max_n\|J_0^{-1}\|$,
$C_J=\max_n\|J_h-J_0\|/h$, observed inverse ratio $\|J_h^{-1}\|/\|J_0^{-1}\|$ and Neumann threshold
$h_0=1/(2M_0C_J)$, in the Euclidean norm on the stage layout (E) and in the endpoint-linearized
norm $\|J_0^{-1}\,\cdot\,\|$ (L, where $M_0=1$); the last two columns repeat the L-norm quantities
for the same mechanism without joint friction.}
\label{tab:p2-constants-check}
\scriptsize
\setlength{\tabcolsep}{3pt}
\begin{tabular}{@{}lrlrllllr@{}}
\toprule
$h$ & $M_0$ (E) & $C_J$ (E) & ratio & $h_0$ (E) & $C_J$ (L) & $h_0$ (L) & $C_J$ (L, no fr.) & $h_0$ (L, no fr.) \\
\midrule
0.04 & 18.2 & $7.5\mathrm{e}{2}$ & 1.55 & $3.7\mathrm{e}{-5}$ & $5.1\mathrm{e}{2}$ & $9.8\mathrm{e}{-4}$ & 34.1 & $1.5\mathrm{e}{-2}$ \\
0.02 & 24.3 & $8.2\mathrm{e}{2}$ & 1.26 & $2.5\mathrm{e}{-5}$ & $5.6\mathrm{e}{2}$ & $8.9\mathrm{e}{-4}$ & 35.5 & $1.4\mathrm{e}{-2}$ \\
0.01 & 28.4 & $9.6\mathrm{e}{2}$ & 1.18 & $1.8\mathrm{e}{-5}$ & $6.5\mathrm{e}{2}$ & $7.6\mathrm{e}{-4}$ & 36.5 & $1.4\mathrm{e}{-2}$ \\
\bottomrule
\end{tabular}
\end{table}

\begin{table}[!htb]
\centering
\caption{Predictor distances for Lemma~\ref{lem:newton-envelope} on the chain (Euclidean norm on
the stage layout): distance of the lifted endpoint predictor $Z^{(0)}_\ast$ and of the implemented
predictor to $Z_G$, and the ratio $\|Z^{(1)}-Z_G\|/\|Z^{(0)}-Z_G\|$ of the first full Newton step
from the implemented predictor.}
\label{tab:predictor-check}
\small
\begin{tabular}{P{0.08\linewidth}P{0.20\linewidth}P{0.22\linewidth}P{0.20\linewidth}}
\toprule
$h$ & lifted predictor & implemented predictor & first-step ratio \\
\midrule
0.04 & 3.23 & 50.2 & 1.12 \\
0.02 & 2.11 & 50.3 & 1.16 \\
0.01 & 1.07 & 50.4 & 1.18 \\
\bottomrule
\end{tabular}
\end{table}

\paragraph{Solver envelope}
Table~\ref{tab:e7} records, for every grid of the sweep, the largest residual of the accepted Newton
iterate under the production stopping rule and its ratio to $h^7$, which is the constant $c_\eta$
with which that grid satisfies (H3). The accepted residual is set by the tolerance, not by $h$, so
the ratio grows from $10^{-4}$ to $10^{6}$ across the sweep; as discussed after
Lemma~\ref{lem:newton-envelope}, the resulting bound $2M\|F_h(\widetilde Z)\|$ on the stage
perturbation is pessimistic by three orders of magnitude relative to the observed errors. The
endpoint position-level constraint norm of Table~\ref{tab:e1} is the quantity that actually
dominates the deviation of the numerical trajectory from the constraint manifold; it decreases as
$h^6$ as the theorem requires.

\begin{table}[!htb]
\centering
\caption{Solver envelope on the chain sweep ($\tau_N=10^{-11}$): largest residual norm of the
accepted iterate over all steps, its ratio to $h^7$ (the $c_\eta$ of (H3) for that grid), and the
largest components of the stage displacements $d_j,\dot d_j,\ddot d_j$ orthogonal to $n_j$.}
\label{tab:e7}
\footnotesize
\begin{tabular}{@{}lrlll@{}}
\toprule
$h$ & Newton/step & $\max_n\|F_h(\widetilde Z_n)\|$ & ratio to $h^7$ ($c_\eta$) & max $\perp n_j$ component \\
\midrule
$0.1$      & 6.00 & $8.1\mathrm{e}{-12}$ & $8.1\mathrm{e}{-5}$ & $5.6\mathrm{e}{-13}$ \\
$0.05$     & 6.00 & $3.5\mathrm{e}{-13}$ & $4.5\mathrm{e}{-4}$ & $3.2\mathrm{e}{-15}$ \\
$0.025$    & 5.90 & $8.8\mathrm{e}{-12}$ & $1.4$               & $6.4\mathrm{e}{-14}$ \\
$0.0125$   & 5.72 & $2.7\mathrm{e}{-12}$ & $5.7\mathrm{e}{1}$  & $4.3\mathrm{e}{-15}$ \\
$0.00625$  & 5.72 & $3.1\mathrm{e}{-12}$ & $8.3\mathrm{e}{3}$  & $4.1\mathrm{e}{-15}$ \\
$0.003125$ & 5.72 & $3.4\mathrm{e}{-12}$ & $1.2\mathrm{e}{6}$  & $5.3\mathrm{e}{-15}$ \\
\bottomrule
\end{tabular}
\end{table}

\subsection{The Gauss family: orders two, four and six at the same Newton work}
\label{sec:e2}

Replacing the three-stage tableau by the two-stage Gauss tableau (order four) or by the implicit
midpoint rule (order two) leaves the stage system of Eq.~\eqref{eq:g6fullva-expanded-stage}
unchanged except for its size ($88$ and $44$ unknowns). Figure~\ref{fig:e2} and
Table~\ref{tab:e2} compare the three members on the chain over $[0,0.5]$ against the same
reference. The fitted position orders are $2.10$, $4.17$ and $6.25$, in agreement with the remark
after Theorem~\ref{thm:g6fullva-order}. The Newton iteration count per step is the same for the
three members ($5.7$ to $6.0$), so at a given step size the work differs only by the size of the
dense factorization; in the present Python implementation the wall time per step differs by at most
a third across the family, being dominated by interpreter overhead. At $h=0.0125$ the three-stage
member is $1.6\times10^{5}$ times more accurate than the one-stage and $310$ times more accurate
than the two-stage member, and at $h=0.003125$ the factors are $4\times10^{7}$ and
$5\times10^{3}$; to reach the accuracy of the three-stage member at $h=0.05$ ($2\times10^{-6}$)
the one-stage member needs $h\approx2\times10^{-3}$, that is, about $23$ times more steps and
Newton iterations.

\begin{table}[!htb]
\centering
\caption{Gauss family with the same stage system on the chain, $T=0.5$: position and velocity
errors at $t=T$ and total Newton iterations. The wall time of the first row of each member includes
just-in-time compilation and is omitted.}
\label{tab:e2}
\scriptsize
\setlength{\tabcolsep}{3pt}
\begin{tabular}{@{}lllrlllr@{}}
\toprule
$h$ & \multicolumn{3}{l}{1 stage (order 2)} & \multicolumn{3}{l}{2 stages (order 4)} & \\
    & position & velocity & Newton & position & velocity & Newton & \\
\midrule
$0.1$      & $1.16\mathrm{e}{-2}$ & $6.53\mathrm{e}{-2}$ & 31   & $1.48\mathrm{e}{-3}$ & $3.23\mathrm{e}{-3}$ & 30 & \\
$0.05$     & $1.16\mathrm{e}{-3}$ & $1.21\mathrm{e}{-2}$ & 60   & $6.23\mathrm{e}{-5}$ & $8.38\mathrm{e}{-5}$ & 60 & \\
$0.025$    & $2.95\mathrm{e}{-4}$ & $2.97\mathrm{e}{-3}$ & 115  & $1.61\mathrm{e}{-6}$ & $2.81\mathrm{e}{-6}$ & 117 & \\
$0.0125$   & $7.41\mathrm{e}{-5}$ & $7.39\mathrm{e}{-4}$ & 229  & $1.45\mathrm{e}{-7}$ & $2.66\mathrm{e}{-7}$ & 229 & \\
$0.00625$  & $1.87\mathrm{e}{-5}$ & $1.85\mathrm{e}{-4}$ & 457  & $9.09\mathrm{e}{-9}$ & $1.67\mathrm{e}{-8}$ & 458 & \\
$0.003125$ & $4.68\mathrm{e}{-6}$ & $4.61\mathrm{e}{-5}$ & 915  & $5.69\mathrm{e}{-10}$ & $1.04\mathrm{e}{-9}$ & 915 & \\
$0.0015625$ & $1.17\mathrm{e}{-6}$ & $1.15\mathrm{e}{-5}$ & 1829 & $3.55\mathrm{e}{-11}$ & $6.51\mathrm{e}{-11}$ & 1830 & \\
$0.00078125$ & $2.93\mathrm{e}{-7}$ & $2.88\mathrm{e}{-6}$ & 3658 & -- & -- & -- & \\
\midrule
fitted order & 2.10 & 2.04 & & 4.17 & 4.17 & & \\
\bottomrule
\end{tabular}
\end{table}

\paperfig{\figpath/e2_gauss_family.png}
{Gauss family with the same FullVA stage system on the chain, $T=0.5$: position error against step
size (left) and against total Newton iterations (right). Three stages (order six) are the method of
this paper.}
{fig:e2}

\subsection{Sharpness of the friction law}
\label{sec:e3}

The Brown--McPhee law is smooth for every $v_s>0$, but its derivatives grow like $v_s^{-k}$, and
Lemma~\ref{lem:p2-from-p1} predicts a small-step threshold that shrinks accordingly.
Figure~\ref{fig:e3} and Table~\ref{tab:e3} repeat the convergence sweep of Section~\ref{sec:e1}
for $v_s\in\{0.5,0.2,0.1,0.05,0.02\}$, each against its own reference at $h_{\rm ref}=0.1/256$.
Two things happen as $v_s$ decreases. The fitted order over the whole sweep drops from $6.25$ to
$2.47$, because the coarse grids are in a pre-asymptotic regime in which the error is dominated by
the friction transition; and the step size at which the pairwise order returns to six moves from
$h=0.1$ to $h\approx3\times10^{-3}$, where all five cases show pairwise orders between $5.3$ and
$6.7$ on the finest pair. The measured Neumann threshold $h_0$ decreases from $7.6\times10^{-4}$ to
$7.1\times10^{-6}$ over the same range, so the analysis predicts the trend correctly but its
sufficient condition is two to three orders of magnitude more conservative than the observed
recovery step. Every run in the sweep converged, with $5.7$ to $6.4$ Newton iterations per step.

\begin{table}[!htb]
\centering
\caption{Friction sweep on the chain, $T=0.5$: fitted position and velocity orders over the six
step sizes, pairwise position orders from $h=0.1$ to $h=0.003125$, coarsest step at which the
pairwise order reaches $5.5$, and the Neumann threshold $h_0=1/(2C_J)$ in the endpoint-linearized
norm measured over the first ten steps at $h=0.01$.}
\label{tab:e3}
\scriptsize
\setlength{\tabcolsep}{3.5pt}
\begin{tabular}{@{}lrrlll@{}}
\toprule
$v_s$ & position order & velocity order & pairwise position orders & recovery $h$ & $h_0$ (L) \\
\midrule
0.5  & 6.25 & 6.11 & 6.95, 4.58, 7.64, 5.98, 5.99 & 0.1     & $7.6\mathrm{e}{-4}$ \\
0.2  & 4.84 & 4.74 & 2.11, 4.54, 2.74, 8.57, 5.85 & 0.0125  & $1.4\mathrm{e}{-4}$ \\
0.1  & 4.74 & 4.78 & 0.47, 5.71, 4.14, 5.88, 6.72 & 0.05    & $6.7\mathrm{e}{-5}$ \\
0.05 & 3.66 & 3.66 & 0.17, 4.11, 4.47, 3.46, 5.27 & --      & $2.1\mathrm{e}{-5}$ \\
0.02 & 2.47 & 2.48 & $-0.22$, 2.29, 2.90, 1.55, 6.20 & 0.00625 & $7.1\mathrm{e}{-6}$ \\
\bottomrule
\end{tabular}
\end{table}

\paperfig{\figpath/e3_friction_sweep.png}
{Friction sweep on the chain: position error against step size for five Stribeck velocities (left)
and fitted orders against $v_s$ (right).}
{fig:e3}

\subsection{The four benchmark mechanisms}
\label{sec:e4}

The public benchmark suite of~\cite{kissel2022absolute} consists of a driven single pendulum, a
free double pendulum released from the horizontal, and two driven closed loops, a four-link
mechanism and a slider-crank (\ref{app:lower-pair}). Three of the four are kinematically
driven: their motion is fixed by the drive constraint and the dynamics only determines the
reactions. \method{} is run on all four over $T=1$ with $h=0.1/2^k$ (Figure~\ref{fig:e4},
Table~\ref{tab:e4}). For the driven single pendulum, whose exact motion is known in closed form,
the final-state position errors are $5.9\times10^{-10}$, $9.3\times10^{-12}$ and
$1.5\times10^{-13}$ at $h=0.1$, $0.05$ and $0.025$ (pairwise orders $6.0$ and $6.0$) and at
roundoff, $10^{-14}$ to $10^{-15}$, for every finer step; its velocity errors scatter between
$10^{-15}$ and $4\times10^{-9}$ without trend in $h$, at the level set by the Newton tolerance of
that harness, which does not project the endpoint velocities. For the two closed loops the trajectory errors against the
$h=0.1/128$ reference are between $6\times10^{-16}$ and $1.3\times10^{-14}$ at every step size
from $0.05$ down to $0.0016$: the driven motion is reproduced to roundoff, the loop-closure
constraints included, and no order can be measured. The double pendulum is the only genuine
dynamics test of the suite. Its final-state errors against the $h=0.1/128$ reference are
$4.6\times10^{-5}$, $1.6\times10^{-7}$, $1.2\times10^{-9}$, $1.9\times10^{-11}$ and
$3.1\times10^{-13}$ for $h=0.1$ to $0.00625$, with pairwise orders $8.2$, $7.0$, $6.0$, $6.0$ in
position and $9.7$, $7.0$, $6.0$, $5.9$ in velocity; the coarsest pair is pre-asymptotic, the rest
is the sixth order of Theorem~\ref{thm:g6fullva-order}. At $h=0.003125$ the error,
$2.3\times10^{-14}$, has reached the resolution floor of the reference ($2\times10^{-14}$,
the difference between the $h=0.1/128$ and $h=0.1/64$ runs).

\paperfig{\figpath/mechanisms.png}
{The four benchmark mechanisms of the public suite~\cite{kissel2022absolute}; joint data in
\ref{app:lower-pair}.}
{fig:mechanisms}

\begin{table}[!htb]
\centering
\caption{The four benchmark mechanisms over $T=1$: final-state position and velocity errors (single
pendulum: against the analytic driven motion; double pendulum: against the $h=0.1/128$ reference)
and trajectory $L^\infty$ errors of the closed loops against their $h=0.1/128$ reference. The
pairwise order refers to the position error and the next finer step; the double-pendulum reference
has a resolution floor of $2\times10^{-14}$.}
\label{tab:e4}
\scriptsize
\setlength{\tabcolsep}{3pt}
\begin{tabular}{@{}lllrllr@{}}
\toprule
$h$ & \multicolumn{3}{l}{single pendulum (driven, analytic reference)} & \multicolumn{3}{l}{double pendulum (free)} \\
    & position & velocity & order & position & velocity & order \\
\midrule
$0.1$      & $5.94\mathrm{e}{-10}$ & $2.43\mathrm{e}{-9}$  & 6.0 & $4.59\mathrm{e}{-5}$  & $1.56\mathrm{e}{-4}$  & 8.2 \\
$0.05$     & $9.32\mathrm{e}{-12}$ & $3.83\mathrm{e}{-11}$ & 6.0 & $1.58\mathrm{e}{-7}$  & $1.85\mathrm{e}{-7}$  & 7.0 \\
$0.025$    & $1.45\mathrm{e}{-13}$ & $9.49\mathrm{e}{-11}$ & --  & $1.24\mathrm{e}{-9}$  & $1.48\mathrm{e}{-9}$  & 6.0 \\
$0.0125$   & $1.08\mathrm{e}{-14}$ & $6.77\mathrm{e}{-10}$ & --  & $1.93\mathrm{e}{-11}$ & $2.29\mathrm{e}{-11}$ & 6.0 \\
$0.00625$  & $1.58\mathrm{e}{-14}$ & $3.73\mathrm{e}{-9}$  & --  & $3.10\mathrm{e}{-13}$ & $3.86\mathrm{e}{-13}$ & 3.8 \\
$0.003125$ & $1.33\mathrm{e}{-15}$ & $1.05\mathrm{e}{-15}$ & --  & $2.29\mathrm{e}{-14}$ & $6.15\mathrm{e}{-14}$ & -- \\
\midrule
 & \multicolumn{3}{l}{four-link (driven loop)} & \multicolumn{3}{l}{slider-crank (driven loop)} \\
 & position & velocity & & position & velocity & \\
\midrule
$0.05$--$0.0016$ & $0.9$--$1.3\mathrm{e}{-14}$ & $1.5$--$4.0\mathrm{e}{-14}$ & -- & $0.6$--$1.2\mathrm{e}{-15}$ & $1.4$--$2.8\mathrm{e}{-15}$ & -- \\
\bottomrule
\end{tabular}
\end{table}

\paperfig{\figpath/e4_benchmarks.png}
{\method{} on the four benchmark mechanisms over $T=1$. The three driven mechanisms are reproduced at
roundoff for every step size; the double pendulum shows the sixth order.}
{fig:e4}

\subsection{Comparison with the public absolute-coordinate codes}
\label{sec:e5}

The public codes of the benchmark suite (the 2021 rA, rp and r$\varepsilon$ formulations of
\cite{kissel2022absolute} and the 2022 half-implicit rA and rA-half codes of
\cite{fang2023halfimplicit}) are run as published, in their dynamics mode with their own tolerance
policies, on the double pendulum over the public horizon $T=3$. Every error in
Table~\ref{tab:e5} is the final-state $L^\infty$ difference to one common reference, the local
\method{} run at $h=10^{-4}$ ($30\,000$ steps). That the public model and the local one are the same
mechanism is checked on the short horizon $T=0.1$, where the public rA final state differs from
the local one by $3.7\times10^{-3}$ at $h=0.1/16$ and $9.2\times10^{-4}$ at $h=0.1/64$, a
first-order decrease and not a constant offset.

Over $T=3$ the double pendulum amplifies perturbations strongly, and the public codes do not
resolve it at any of the step sizes: the rA, r$\varepsilon$ and 2022 rA codes give identical
results with position errors from $1.8$ at $h=0.1$ to $0.17$ at $h=0.001$ (pairwise orders
$0.2$ to $0.9$, approaching first order only below $h=0.002$); the rp code fails to converge for
$h\ge0.025$ and is at first order below; the half-implicit rA-half code does not converge at all
on this horizon (errors between $0.9$ and $5.8$). \method{} reaches $4.3\times10^{-4}$ at
$h=0.1$, $2.1\times10^{-11}$ at $h=0.01$ and the resolution floor of the reference,
$2\times10^{-13}$, at $h\le0.002$, with pairwise orders $7.2$, $9.2$, $6.0$ and $6.1$ on the first
four pairs. The Newton counts are $3.0$ to $4.5$ iterations per step for \method{} against $1.8$
to $6.8$ for the 2021 codes and $3$ to $32$ for the 2022 codes; wall times are not comparable
across the implementations (the local harness is a Python reference implementation with a dense
automatic-differentiation Jacobian) and are omitted.

On the three driven mechanisms the public horizon adds nothing to Section~\ref{sec:e4}:
\method{} reproduces the driven motion to $10^{-14}$ over $T=3$ at $h=10^{-2}$, $10^{-3}$ and
$10^{-4}$, while the public rA code, measured against its own kinematic-mode reference, has
position errors below $10^{-9}$ but velocity errors that decrease at first order,
$1.1\times10^{-1}$, $1.1\times10^{-2}$, $1.1\times10^{-3}$ on the four-link mechanism and
$8.5\times10^{-3}$, $8.7\times10^{-4}$, $8.8\times10^{-5}$ on the slider-crank.

\begin{table}[!htb]
\centering
\caption{Double pendulum over the public horizon $T=3$: final-state $L^\infty$ errors against the
common \method{} reference at $h=10^{-4}$. The 2021 r$\varepsilon$ and the 2022 half-implicit rA
codes return the same numbers as the 2021 rA code. ``fail'': the public Newton iteration did not
converge. The reference has a resolution floor of about $2\times10^{-13}$.}
\label{tab:e5}
\scriptsize
\setlength{\tabcolsep}{3.5pt}
\begin{tabular}{@{}lllllll@{}}
\toprule
$h$ & \multicolumn{2}{l}{\method{}} & \multicolumn{2}{l}{2021 public rA} & 2021 public rp & 2022 rA-half \\
    & position & velocity & position & velocity & position & position \\
\midrule
$0.1$    & $4.29\mathrm{e}{-4}$  & $1.51\mathrm{e}{-3}$  & $1.84$ & $6.04$ & fail   & $5.76$ \\
$0.05$   & $3.03\mathrm{e}{-6}$  & $9.84\mathrm{e}{-6}$  & $1.58$ & $5.03$ & fail   & $5.01$ \\
$0.025$  & $5.34\mathrm{e}{-9}$  & $1.72\mathrm{e}{-8}$  & $1.34$ & $4.13$ & fail   & $1.34$ \\
$0.0125$ & $8.12\mathrm{e}{-11}$ & $2.66\mathrm{e}{-10}$ & $1.06$ & $3.43$ & $0.749$ & $0.903$ \\
$0.01$   & $2.10\mathrm{e}{-11}$ & $7.09\mathrm{e}{-11}$ & $0.958$ & $3.40$ & $1.25$ & $1.12$ \\
$0.002$  & $2.56\mathrm{e}{-13}$ & $1.76\mathrm{e}{-12}$ & $0.319$ & $1.10$ & $0.233$ & $4.26$ \\
$0.001$  & $2.16\mathrm{e}{-13}$ & $1.66\mathrm{e}{-12}$ & $0.169$ & $0.532$ & $0.113$ & $1.92$ \\
\bottomrule
\end{tabular}
\end{table}

\paperfig{\figpath/e5_double_pendulum.png}
{Double pendulum over the public horizon $T=3$: final position (left) and velocity (right) errors
of \method{} and of the public absolute-coordinate codes against the common $h=10^{-4}$ reference.}
{fig:e5}


\subsection{Constraint drift and energy over the regular branch}
\label{sec:e6}

Figure~\ref{fig:e6} follows the chain over the whole regular branch $[0,0.5]$ at $h=0.01$ in two
variants: the frictional benchmark, and a conservative variant with $\mu_s=\mu_d=0$, no viscous
damping and no external loads, for which the total mechanical energy is an invariant of the exact
flow. The endpoint position- and velocity-level constraint norms stay at $1.1\times10^{-13}$ and
$9.5\times10^{-14}$ in both variants, the acceleration-level norms at the stages below
$2\times10^{-14}$, and the quaternion unit error at $2\times10^{-16}$; the relative energy error of
the conservative variant is $4.3\times10^{-12}$ at every step (the exact Gauss step is symplectic
on the reduced problem, and the endpoint closure perturbs it by $O(h^7)$), while the frictional
variant dissipates $37\%$ of its initial energy over the branch. The Newton count is $5$ per step
without friction and $5.7$ with it.

\paperfig{\figpath/e6_regular_branch.png}
{Chain over the regular branch $[0,0.5]$ at $h=0.01$: relative energy error of the conservative
variant (left) and endpoint constraint norms of both variants (right; the velocity-level closure is
applied only when the raw defect exceeds $10^{-13}$, hence the saw-tooth).}
{fig:e6}
